% Source PDF: source/普通高考/2007/2007山东理.pdf, page 1, question 10.
% Node→directed-edge list:
% 开始→输入 n→S=0,T=0→n<2?；是→输出 S,T→结束；
% 否→S=S+n→n=n-1→T=T+n→n=n-1→回到判定框前共享入口。
\begin{tikzpicture}[
    >=Stealth,
    line width=0.45pt,
    line cap=round,
    line join=round,
    every node/.style={
        anchor=center,
        align=center,
        font=\normalsize,
        execute at begin node={\setlength{\parindent}{0pt}}
    },
    flowbox/.style={
        draw,
        inner xsep=3pt,
        inner ysep=1.5pt,
        minimum height=0.65cm,
        align=center,
        execute at begin node={\setlength{\parindent}{0pt}}
    },
    terminal/.style={flowbox,rounded corners=2.5pt,minimum width=1.15cm},
    io/.style={flowbox,shape=trapezium,trapezium left angle=72,trapezium right angle=108},
    decision/.style={flowbox,shape=diamond,aspect=2.7,inner sep=1pt,minimum height=0.9cm},
    flowarrow/.style={->}
]
    \node[terminal] (start) at (0,0) {开始};
    \node[io,minimum width=1.55cm] (input) at (0,-1.2) {输入 $n$};
    \node[flowbox,minimum width=2.6cm] (init) at (0,-2.4) {$S=0,T=0$};
    \coordinate (join) at (0,-3.05);
    \node[decision,minimum width=2.7cm] (test) at (0,-3.8) {$n<2?$};
    \node[flowbox,minimum width=2.0cm] (sum) at (0,-5.05) {$S=S+n$};
    \node[flowbox,minimum width=2.0cm] (dec1) at (0,-6.25) {$n=n-1$};
    \node[flowbox,minimum width=2.2cm] (addt) at (0,-7.45) {$T=T+n$};
    \node[flowbox,minimum width=2.0cm] (dec2) at (0,-8.65) {$n=n-1$};
    \node[io,minimum width=2.25cm] (output) at (3.75,-5.05) {输出 $S,T$};
    \node[terminal] (finish) at (3.75,-6.35) {结束};

    \draw (start.south) -- (input.north);
    \draw (input.south) -- (init.north);
    \draw (init.south) -- (join);
    \draw[flowarrow] (join) -- (test.north);
    \draw[flowarrow] (test.south) -- node[right=2pt] {否} (sum.north);
    \draw[flowarrow] (sum.south) -- (dec1.north);
    \draw[flowarrow] (dec1.south) -- (addt.north);
    \draw[flowarrow] (addt.south) -- (dec2.north);
    \coordinate (exit-right) at (3.75,-3.8);
    \draw (test.east) -- node[above=2pt] {是} (exit-right);
    \draw[flowarrow] (exit-right) -- (output.north);
    \draw[flowarrow] (output.south) -- (finish.north);

    \coordinate (loop-bottom) at (-2.25,-9.05);
    \coordinate (loop-top) at (-2.25,-3.05);
    \draw (dec2.south) -- (0,-9.05) -- (loop-bottom) -- (loop-top);
    \draw[flowarrow] (loop-top) -- (join);
\end{tikzpicture}
